\section{Operador Nabla ($\nabla$)}
$\nabla$ é um operador que age em funções diferenciáveis cujo domínio são vetores em $\mathbb{R}^n$. Em coordenadas cartesianas $(x^1, \dots , x^n)$ e base natural $\{\mathbf{e}_1, \dots, \mathbf{e}_n\}$, podemos considerá-lo como sendo um "vetor de operadores derivada":  $\nabla \equiv \mathbf{e}_1\partial_1 + \dots + \mathbf{e}_n\partial_n$. Por si só, essa expressão não tem significado, e sua interpretação depende da natureza função.

Se $F$ é um campo escalar, $\nabla F$ é o \textit{gradiente} de $F$. Considere um deslocamento infinitesimal $d\mathbf{x} = \mathbf{u}dt$, onde $\mathbf{u}$ é um vetor unitário, de modo que $dx^i = u^i dt$. A variação $dF$ é dada por
\begin{equation*}
F(\mathbf{x}+ d\mathbf{x}) - F(\mathbf{x}) \equiv dF = \partial_i F \, dx^i = dt(\mathbf u \cdot \nabla F),
\end{equation*}
\noindent portanto, $\nabla F$ aponta na direção de maior taxa de variação de $F$, e é normal à superfícies $F =$ constante. Na equação acima e no resto do trabalho, foi utilizada a notação de Eistein, onde índices repetidos em posições diferentes (em cima e em baixo) significam que há um somatório implícito sobre esse índice.

Se $\mathbf{F} = F^i \mathbf{e}_i$ é um campo vetorial, $\nabla \cdot \mathbf{F}$ é a \textit{divergência} de $\mathbf{F}
$, e $\nabla \times \mathbf{F}$ é seu \textit{rotacional}. A divergência em um ponto $\mathbf{x}$ é definida como
\begin{equation}
\nabla \cdot \mathbf{F} \equiv \lim\limits_{V \to 0} \frac{\Phi_F(S)}{V} ,
\label{def: div}
\end{equation}

\noindent onde $V$ é o volume da região delimitada por $S$, e $\mathbf{x}$ é um ponto dentro dessa região. Se $\nabla \cdot \mathbf{F} > 0$, então $\mathbf{x}$ é uma fonte de $\mathbf{F}$; se $\nabla \cdot \mathbf{F} < 0$, então $\mathbf{x}$ é um sorvedouro de $\mathbf{F}$. Em coordenadas cartesianas, a expressão para a divergência de um campo vetorial é $\nabla \cdot \mathbf{F} = \partial_i F^i$.

O rotacional de um campo vetorial, diferentemente das outras aplicações do operador $\nabla$, requer que $\mathbf{F}$ seja uma função $\mathbb{R}^3 \to \mathbb{R}^3$, sendo $\nabla \times \mathbf{F}$ também uma função desse tipo. A componente de $\nabla \times \mathbf{F}$ na direção de um vetor unitário $\mathbf{n}$ num ponto $\mathbf{P}$ é definida como
\begin{equation}
\mathbf{n} \cdot \nabla \times \mathbf{F} \equiv \lim\limits_{A \to 0} \frac{1}{A} \oint_C \mathbf{F} \cdot d\mathbf{l},
\label{def: rot}
\end{equation}

\noindent onde $A$ é a área de uma região do plano perpendicular a $\mathbf{n}$ que passa por $\mathbf{P}$ e contém esse ponto, e $C$ é um circuito fechado que delimita essa área. Em coordenadas cartesianas, a fórmula para o rotacional é 
\begin{equation}
\nabla \times \mathbf{F} = \epsilon^{ijk}\partial_i F_j \mathbf{e}_k,
\label{eq: xyz_rot}
\end{equation}

\noindent onde $\epsilon^{ijk}$ é o tensor de Levi-Civita.
